% k02 King's existence theorem on the cube with scalar sources
\section{King's existence theorem on the cube with scalar sources}\label{sec:k02}

This section proves \cref{hyp:s03:W1}, the existence theorem with scalar sources in the form the Wick leg consumes: in the reference scheme $z$, the fixed-cube limits of the generating functional with curvature and scalar sources exist, locally uniformly and holomorphically, on the domain of coefficient vectors whose positive scalar part is below a threshold, and they are positive at real points; the negative scalar part is unrestricted. The route is the Cauchy argument of \cref{sec:k01} with the scheme $c$-number carried along (\cref{thm:k02:CAh}), applied at real points with a small non-negative scalar source, followed by a normal-family argument on the domain. The domain is bounded by Hölder's inequality from three one-sided estimates: the curvature bound of \cref{lem:k01:free}, a positive-source bound (\cref{prop:k02:Wsharp}), and a negative-source bound (\cref{thm:k02:UB}), the last obtained from positive-source statements at a heavier scalar mass by an exact mass shift and a correlation inequality. The notation is that of \cref{ssec:k01:setting}.

\subsection{Setting; the source-dependent graphs}\label{ssec:k02:setting}

Throughout, the scheme is the reference scheme $z=a^\sharp$ of \cref{def:s02:sources}; write $z_k$ for its $k$-th term. Put $\sS_k(h)\defeq\sum_{x\in\Lk\cap\CJ}\eps_k^3h(x)$ and $c_k(h)\defeq z_k\sS_k(h)$, which is linear in $h$. For a real source $s=(u,h)$ the normalised functional in the scheme is
\[
n_k(s)\ \defeq\ e^{-c_k(h)}\,N_k(s)\ =\ e^{-z_k\sS_k(h)}\,\avK{e^{F(u)+\Wick{h}}}\ =\ Z^{z,K}_{k,J}(s),
\]
the object of \cref{hyp:s03:W1}. For a finite family $f_1,\dots,f_n$ of real $2$-forms and $h_1,\dots,h_p$ of real functions, all supported at $\ell^\infty$-distance $>1$ from $\R^3\setminus\CJ$, and $\zeta=(z,w)\in\C^{n+p}$, $n_k(\zeta)\defeq e^{-\sum_jw_jc_k(h_j)}\avK{e^{\sum_iz_iF(f_i)+\sum_jw_j\Wick{h_j}}}$ (here $z_i$ are coefficients and $z_k$ the scheme term; no confusion will arise). Every $n_k$ is entire in $\zeta$, $|n_k(\zeta)|\le n_k(\Re\zeta)$, and $\ln n_k$ is convex on real points (\cref{lem:s02:alg}(d)): a log-Laplace transform minus a linear function.

A second scalar mass $M>0$ will be used with the \emph{same} $e$, $\lambda$, $\lmu$, lattice, cube and boundary conditions; $\langle\cdot\rangle^{(M)}$ denotes King's cube state at mass $M$, $\alpha_k(M)\defeq M^2+\dm_{\eps_k}(M)$ King's free-boundary bare coefficient at mass $M$ (as in \cref{hyp:s03:T2}), $u_k(M)\defeq\frac12(\alpha_k(M)-\alpha_k(m))$, and $s_-(M)\defeq\liminf_ku_k(M)$. Every statement of \cref{sec:k01} and of this section holds at every point of King's range with point-dependent constants, in particular at the point $(e,\lambda,M,\lmu)$; we use this at $M=m$ and at one heavier $M_2$.

\begin{lemma}[\Kh: one-step convergence of the source-dependent graphs on King's cube]\label{hyp:k02:Kh}
Let $P^{(j),1}_k$ be the interaction part of King's effective action after $j$ steps on $\CK$ \cite[(3.17)]{K1}, written as a sum over connected renormalised graphs $H$, and let $\mathcal N_h$ be the family of its members that depend on the scalar source $h$: the field-dependent part of every graph with an $h$-vertex, and every field-independent $h$-dependent graph except the $h$-linear ones of the three lowest orders $(\alpha,\beta)\in\{(0,0),(0,1),(2,0)\}$ in $(e^2,\lambda)$. Then for every fixed real $h\in C^\infty_c$ supported at distance $>1$ from $\R^3\setminus\CJ$ and every member of $\mathcal N_h$, for $J_*\le j\le k-J_*$: the localisation bound \cite[(3.39)]{K1} holds for its localisations at distance from the faces larger than the small-field scale; the bound \cite[(3.38)]{K1} with the restricted degree holds with the factor $(1+k-j)^{p_0}$, summed over the localisations sharing the cube of an $h$-vertex, and for each single localisation at distance $>1$ from the faces; and the sum over positions of the difference between the member at $(k,j)$ and the member at $(k+1,j+1)$, on the common unit lattice, is bounded by $C\big[5^{-\gamma'j}\eta^{-\beta'}+\eta^{\sigma'}\big]\,|\CJ|$ with $C$ independent of $j$ and $k$.
\end{lemma}
This is the cube form, on $\CK$ with King's Dirichlet weight, of the corresponding statement on King's tori; the $h$-free graphs, those with a curvature vertex and no $h$-vertex and the source-free ones, are covered on the cube by \cref{prop:k01:Kpos}. The localisations of a member of $\mathcal N_h$ with tree length $R\le p$ lie at distance $>1$ from the faces, where the bounds reduce to those of the free flow on the torus (\cref{prop:k02:KhI}); a localisation with $R>p$ carries the factor $e^{-\delta R}$ once the localised graph is cut into clusters of mutual distance $>3\sqrt d$, each reduced in its own box of side $O(p)$, the clusters containing the source vertex lying away from the faces and the clusters meeting the Dirichlet faces being source-free (\cref{prop:k02:KhR}). The proof is given in \cref{ssec:k02:Khproof}.
% k0x_fill_Kh: proof of hyp:k02:Kh on King's cube, both branches (input from sec/k02_sources.tex after the lemma's commentary).
% Uses the units, zones and the cut by clusters of the appendix on prop:k01:Kpos (ssec:a02:zones, a02_deg2); no macro of that appendix is used here.

\subsection{Proof of \texorpdfstring{\cref{hyp:k02:Kh}}{the one-step convergence of the source-dependent graphs}}\label{ssec:k02:Khproof}

We work in the unit units of depth $j$ of \cref{ssec:a02:zones}: the live sites are $T_\eta$, $F$ denotes the scalar Dirichlet planes, $t(x)\defeq\dist(x,F)$, and a localised graph $H$ carries King's boxes $\square\subset\square''\subset\square'$ and $\Omega=\square'\cap T_\eta$, in zone~I if $\square'\cap F=\emptyset$ and in zone~B otherwise. The localisation is by centres, as in the proof of \cref{prop:k01:Kpos}: a tuple $\tau$ has the point set $P(\tau)$, $N\defeq|P(\tau)|\le N_*$ with $N_*$ depending on the order $\bar n$ of the expansion only, and the tree length $R=R(\tau)\defeq\dist(P(\tau))$. The cut by clusters is that of \cref{lem:a02:CL}: $c_0\defeq\frac{13}3\sqrt3$, $R_0\defeq3\sqrt3$, the clusters $K$ are the components of the $c_0$-proximity graph on $P(\tau)$, every line whose ends lie in different clusters is cut and kept whole as an external line of both cluster graphs $H_K$, and $H_K$ is reduced with its own boxes $\square_K\subset\square''_K\subset\square'_K$ and gauge $B_K=A_j(x^K_0)$; put $\Omega_K\defeq\square'_K\cap T_\eta$. A cluster is in zone~I if $\Omega_K\cap F=\emptyset$ (then $\square'_K$ is not truncated and $\Omega_K=\square'_K$) and in zone~B otherwise. The two members of the diagonal pair, the $k$-model at depth $j$ and the $(k+1)$-model at depth $j+1$, live on the same unit lattice; every bound below holds for each of them separately with the same constants, and the two are compared only in zone~I.

The $h$-vertex is King's vertex \cite[(3.21)]{K1}: the smooth function $h$ sampled at the points of the lattice's own frame, the coupling $\theta\eta^{2}$ in the units of depth $j$, and at most two fields $\phi^{(j)}$. In King's count of the coupling power $\rho(H)=\sum_v\rho_v$ of a graph \cite[p.~665]{K1}, every vertex has $\rho_v-s_v/4\ge0$, $s_v$ its number of external scalar fields, and the $h$-vertex has $\rho_h-s_h/4\ge2-\frac12=\frac32$; hence
\begin{equation}\label{eq:k02:margin}
\rho(H)-\tfrac14s(H)\ \ge\ \tfrac32\qquad\text{for every member of $\mathcal N_h$ with external fields,}
\end{equation}
$s(H)$ its number of external scalar legs, and $\rho(H)\ge0$ for the field-independent members. The Wick constant of \cite[(3.21)]{K1} is field-free and $h$-linear of the lowest order, so it enters only the excluded family $\mathcal T_h$ of \cref{app:a04:def:families} and no member of $\mathcal N_h$ depends on it.

\paragraph{The range.} Let $\eta_2\le\eta_1$ be so small that for $\eta\le\eta_2$, with $p=p(\eta)$ and $M$ the block integer of $\CJ$,
\begin{equation}\label{eq:k02:range}
\eta^{-1}\ \ge\ 4p+8M+2c_0N_*^2+8,\qquad p\ \ge\ c_0N_*^2,\qquad C_\flat\,e\,\eta^{1/2}\,p^{2}\ \le\ \tfrac12,
\end{equation}
where $C_\flat$ is the constant of \cref{lem:k02:Zh}(e) below times the constant of the difference bound \cite[(3.45)]{K1}. All three hold for $\eta$ small because $p(\eta)$ grows slower than every power of $\eta^{-1}$ and faster than every logarithm. The first condition contains the one of \cref{ssec:a02:zones}, so $J_*$ is taken with $5^{-J_*}\le\eta_2$.

\begin{lemma}[the source vertices and the clusters]\label{lem:k02:Zh}
Let $\eta\le\eta_2$, let $H$ be a member of $\mathcal N_h$, $s$ an $h$-vertex of $H$ with centre $z_s$, and $\tau$ a tuple of $H$. Then: \textup{(a)} $t(z_s)\ge\eta^{-1}-3$. \textup{(b)} If $\tau$ is in zone~B, then $R(\tau)\ge\eta^{-1}/2$. \textup{(c)} If $\tau$ is in zone~B, then $\tau$ has at least two clusters. \textup{(d)} Every cluster $K$ containing an $h$-vertex has $\dist(\square_K,F)\ge2p+M+1$, $\Omega_K=\square'_K$ and $\dist(\square_K,\partial\Omega_K)\ge2p$: it is a zone-I cluster. Hence every zone-B cluster is source-free. \textup{(e)} Every $\square'_K$ has side $\le c_0N_*+4p+6M\le C_{\mathrm{box}}\,p$. The same holds with a curvature vertex in place of the $h$-vertex.
\end{lemma}
\begin{proof}
(a) $h(x)\ne0$ forces $\ell^\infty$-$\dist(x,\R^3\setminus\CJ)>1$ in physical units, that is $>\eta^{-1}$ unit lengths; the planes $F$ lie within one fine spacing $5^{-j}\le1$ of $\partial\CJ$, and the vertex lies in $\supp\chi_{z_s}$, within $\frac23\sqrt3<2$ of its centre. (b) A tree of total length $R$ joins the points of $P(\tau)$, so all of them lie within $R$ of $z_s$; $\square$ lies in the sup-norm cube of half-side $R+M$ about $z_s$, the large blocks within $2p$ of it in that of half-side $R+2p+2M$, and $\square'$ in that of half-side $R+2p+3M$. If $\square'$ meets $F$, then $t(z_s)\le R+2p+3M$, and with (a) and \eqref{eq:k02:range} $R\ge\eta^{-1}-2p-3M-3\ge\eta^{-1}/2$. (c) With one cluster $\operatorname{diam}P(\tau)\le c_0(N-1)$ (\cref{lem:a02:CL}(c)), so $R(\tau)\le(N-1)\operatorname{diam}P(\tau)\le c_0(N_*-1)^2<\eta^{-1}/2$ by \eqref{eq:k02:range}, against (b). (d) If $s\in K$, every point of $P_K$ lies within $c_0(N-1)$ of $z_s$, so $\square_K$ lies in the cube of half-side $c_0N+M$ about $z_s$ and $\square'_K$ in that of half-side $c_0N+2p+3M$. By (a) and \eqref{eq:k02:range}, $t(z_s)-(c_0N+2p+3M)\ge\eta^{-1}-3-c_0N-2p-3M>2p+M+1$. So $\square'_K$ misses $F$, it is not truncated, it contains every large block within $2p$ of $\square_K$, and $\dist(\square_K,F)\ge t(z_s)-c_0N-M\ge2p+M+1$. (e) $\square_K$ has side $\le\operatorname{diam}P_K+2M\le c_0N_*+2M$ and $\square'_K$ adds $2p+2M$ per side; truncation at $\partial T_\eta$ only shrinks it. The curvature source $f$ is supported at the same distance from $\R^3\setminus\CJ$, so the same arithmetic applies to its vertex.
\end{proof}

\begin{lemma}[cut lines vanish on the Dirichlet planes]\label{lem:k02:face}
Let $\ell$ be a scalar, dressed or external cut line of \cref{lem:a02:CL} with its end $u$ in a cluster $K$, in either model, at the background $A^{(j)}$ or at zero background. \textup{(a)} For every evaluation point $u'$ of the $u$-end with $t(u')\le\frac14$ and every operation $O_v$ at the far end (a difference or H\"older quotient of order $\le1+\alpha$, or the identity) whose evaluation set $E_v$ lies at distance $\ge\frac12$ from $u'$,
\[
|O_v\ell(u',\cdot)|\ \le\ C\min(1,t(u'))\,e^{-\delta\dist(u',E_v)} .
\]
\textup{(b)} For a scalar or dressed cut line and evaluation points $u',v'$ of its two ends with $|u'-v'|\ge1$ and $t(u'),t(v')\le\frac14$,
\[
|\ell(u',v')|\ \le\ C\min(1,t(u'))\min(1,t(v'))\,e^{-\delta|u'-v'|} .
\]
The constants are those of \cref{lem:a02:XL}(b).
\end{lemma}
\begin{proof}
Each of the three kinds of line vanishes when its $u$-end lies on $F$: a scalar line by the Dirichlet condition of the cube propagator, a dressed line and an external line through their first factor, which is such a propagator. Let $u'_F\in F$ be the foot of $u'$ and telescope along the $t(u')/5^{-j}$ bonds of the normal segment from $u'_F$ to $u'$; the parallel transporters along a bond are unimodular, so $|O_v\ell(u',\cdot)|\le t(u')\max_b|O_vD_b\ell(\cdot,\cdot)|$, the maximum over the bond differences $D_b$ on that segment. Every such bond lies within $\frac14$ of $u'$, hence at distance $\ge\frac14$ from $E_v$, and the mixed bound of \cref{lem:a02:XL}(b) ($C^{1,\alpha}\otimes C^{1,\alpha}$ at unit scale in both arguments, decay $e^{-\delta_1|u-v|}$, no distance condition) bounds each by $Ce^{-\delta\dist(u',E_v)}$. This is (a). For (b), $\ell(u',v'_F)=0$ by the Dirichlet condition in the second variable (for a dressed line through its last factor); telescoping at $v'$ along its normal segment gives $|\ell(u',v')|\le t(v')\max_b|D_b\ell(u',\cdot)|$ over bonds $b$ within $\frac14$ of $v'$, each at distance $\ge1-\frac14-5^{-j}\ge\frac12$ from $u'$, and (a) with $O_v\defeq D_b$ bounds each by $C\min(1,t(u'))e^{-\delta|u'-v'|}$. Where $t>\frac14$ at an end, $\min(1,t)\ge\frac14$ and \cref{lem:a02:XL}(b) alone gives the bound.
\end{proof}
% k0x_fill_Kh2: the two branches of the proof of hyp:k02:Kh and its assembly (input from sec/k0x_fill_Kh.tex).

\begin{proposition}[the branch $R\le p$: zone~I]\label{prop:k02:KhI}
Let $H$ be a member of $\mathcal N_h$ and $\tau$ a localisation of $H$ in zone~I; by \cref{lem:a02:Z}(i) this covers every localisation with $R\le p$. Then, for each model, on $\supp\chi_j$,
\begin{equation}\label{eq:k02:I1}
|E^{(j)}(H;\tau)|\ \le\ C(1+k-j)^{p_0}\,\eta^{\rho(H)}\,p^{n_1}\,e^{-\delta R},
\end{equation}
which is \cite[(3.39)]{K1}, and across the diagonal pair, on the common unit lattice,
\begin{equation}\label{eq:k02:I2}
|E'^{(j+1)}(H;\tau)-E^{(j)}(H;\tau)|\ \le\ C\,\eta^{\rho(H)-1/2}\,p^{n_1}\big[5^{-\gamma j}+\eta^{3+\sigma+1/2}\big]\,e^{-\delta R}.
\end{equation}
The constants depend on the point of King's range, $J$, $\bar n$, $\norm h_{C^1}$ and, for members with curvature vertices, $\max\norm{f_{\mu\nu}}_{C^{2,\alpha}}$; not on $j$, $k$, $\tau$.
\end{proposition}
\begin{proof}
By \cref{lem:a02:I}, which uses only that the localisation is in zone~I, $E(H)=E(H^{\mathrm{free}})+\mathcal R$ with $H^{\mathrm{free}}$ the same graph on the free $f_D$ flow, every term of $\mathcal R$ carrying an image line of any prescribed degree with the factor $e^{-\delta_0p/2}\le C\eta^{3+\sigma+1/2}$, and the free-boundary counterterm renormalising $H^{\mathrm{free}}$ exactly. For $H^{\mathrm{free}}$ we run King's proof of \cite[Thm.~3.5, Prop.~3.6]{K1} as in the proof of \cref{thm:a02:B4}, with the following points for the $h$-vertex. \emph{Census:} every connected component of an ordered prefix that contains the $h$-vertex has positive degree, the vertex \cite[(3.21)]{K1} carrying the coupling power $2$ against at most two fields, except the field-independent $h$-linear graphs of the three lowest orders $P_0$ \cite[Thm.~3.5]{K1}, which form the family $\mathcal T_h$ of \cref{app:a04:def:families} and are not in $\mathcal N_h$; the $h$-free components obey the degree count of \cref{thm:a02:T}, with the Ward rows by \cref{thm:a02:WR}, and the components with block points are convergent by \cref{lem:a02:BLK}. \emph{Bound:} the ordered multiscale bound \cref{thm:a02:S3}(a) with the root held at the $h$-vertex gives \eqref{eq:k02:I1}; the power of $\eta$ is $\rho(H)$ by \cite[p.~665]{K1}, the fields on $\supp\chi_j$ being bounded by \cite[(3.2)]{K1}, which is the content of \eqref{eq:k02:margin}, and the logarithm comes from the infrared sums as in \cref{thm:a02:B4}. \emph{Comparison:} \cref{thm:a02:S3}(b), one factor at a time \cite[pp.~663--665]{K1}: the lines by \cref{prop:a02:FL}, the $h$-vertex by the sampling of the same smooth $h$ at points of the two frames at distance $\le\sqrt3\,\eps_k$, whose difference is $\le C\eps_k\norm h_{C^1}\le C5^{-\gamma j}$, and the finer slices through the finest index sum; each error term is the graph with one element weakened, which costs the factor $\eta^{-1/2}$ in \eqref{eq:k02:I2}. The identities \cite[(3.43)--(3.55)]{K1} act on the $h$-vertex only through the phases at its scalar ends, which the gauge transformation \cite[(3.46)]{K1} carries. Members with curvature vertices: \cref{lem:a02:SBP}(b)--(d) as in \cref{thm:a02:B4}. The $h$-vertex is bounded by $\norm h_\infty$, and only through it does $h$ enter the constants of \eqref{eq:k02:I1}.
\end{proof}

\begin{proposition}[the branch $R>p$: the cut by clusters]\label{prop:k02:KhR}
Let $\eta\le\eta_2$, let $H$ be a member of $\mathcal N_h$ and $\tau$ a tuple of $H$ with at least two clusters. By \cref{lem:k02:Zh}(c) this covers every zone-B tuple, and by \cref{lem:a02:CL}(c) and \eqref{eq:k02:range} every tuple with $R>p$. Then, for each model separately, on $\supp\chi_j$,
\begin{equation}\label{eq:k02:R1}
|E^{(j)}_N(H;\tau)|\ \le\ C(1+k-j)^{p_0}\,\eta^{\rho(H)}\,p^{n_1}\,e^{-\delta\,\dist(P(\tau))},
\end{equation}
$E_N$ the field-dependent part (the whole value for a field-independent member); this is \cite[(3.38)]{K1} at the tuple. Moreover, for every unit cube $y_*$ and $0<\delta'\le\delta/2$,
\begin{equation}\label{eq:k02:R2}
\sum_{\substack{\tau\ \text{in zone B}\\ \text{with an $h$-vertex in }y_*}}|E^{(j)}_N(H;\tau)|\,e^{\delta'R(\tau)}\ \le\ C(1+k-j)^{p_0}\,\eta^{\rho(H)}\,p^{n_1}\,e^{-\delta/(8\eta)} .
\end{equation}
The constants depend on the point of King's range, $J$, $\bar n$, $\norm h_\infty$ and, for members with curvature vertices, $\max\norm{f_{\mu\nu}}_{C^{2,\alpha}}$.
\end{proposition}
\begin{proof}
At a tuple with at least two clusters every connected component of every cluster graph carries a cut-line leg (\cref{cor:a02:DEGall}(c)), so the vacuum members and the component-wise cancellation of \cite[p.~667]{K1} do not arise, and $E(H;A^{(j)};\tau)$ and $E(H;0;\tau)$ are bounded separately; at zero background the cut lines are the cube propagators at $A=0$ and $B_K=0$. Fix one of the two terms.

\emph{Step 0: regularised cut lines.} Replace every cut line $\ell$ by $\ell^\flat$; by \cref{lem:a02:XL}(a) this changes nothing.

\emph{Step 1: reduction per cluster.} Each $H_K$ is reduced by \cite[(3.43)--(3.55)]{K1} with its own box $\Omega_K$, centre $x_0^K$, gauge $B_K$ and field $\tilde A_K=A^{(j)}-B_K$ on $\Omega_K$ (\cref{lem:a02:CL}), its cut lines keeping the background and receiving the unitary factors \cite[(3.46)]{K1} at their ends. The expansion of the reduced propagator in $\tilde A_K$ converges on $\Omega_K$: $e_j\sup_{\Omega_K}|\tilde A_K|\le Ce_jp\operatorname{diam}\Omega_K\le C_\flat e\eta^{1/2}p^2\le\frac12$ by the difference bound \cite[(3.45)]{K1}, \cref{lem:k02:Zh}(e) and \eqref{eq:k02:range}: this is the reduction of the branch $R\le p$ on a box of side $O(p)$, run inside each cluster. By \cref{lem:a02:CL}(a) the partition, the boxes, the gauges and the cut status of every line are the same for every member of every family met in the nesting and for every position, so renormalisation families are formed inside each $H_K$. A zone-I cluster, in particular every cluster containing a source vertex (\cref{lem:k02:Zh}(d)), satisfies $\dist(\square_K,\partial\Omega_K)\ge2p$ and $\Omega_K\cap F=\emptyset$, so \cref{lem:a02:I} applies to $H_K$: it is its free $f_D$-flow version plus terms carrying an image line of any prescribed degree, with the cut lines as external legs; the vertices of the expansion in $\tilde A_K$ that lie in $\square'_K\setminus\square_K$ are joined to a localised vertex of $K$ by a path of lines one of which is long, and that line supplies a credit of any order, as in step (2) of the proof of \cref{thm:a02:WR} run inside $H_K$ (\cref{thm:a02:WRall}). A zone-B cluster is source-free and is rewritten by \cref{def:a02:RC} within $\Omega_K$.

\emph{Step 2: the cut-line factors.} Every evaluation point of a cut-line end in an output term is joined to a localised anchor of its own cluster by at most $C(\bar n)$ lines of that cluster's rewritten units (\cref{thm:a02:WRall}), and by \cref{lem:a02:XL}(b) the phase-dressed cut line satisfies $|O_uO_v\ell^{\flat,\mathrm{ph}}|\le Ce^{-\delta_2d(K,K')}\times$(transfer factors on the lines of those paths), for all operations $O_u,O_v$ of order $\le1+\alpha$ at its ends, and $\le C$ crudely. The quotient graph (clusters as nodes, cut lines as edges) is connected because $H$ is (\cref{lem:a02:CL}(d)); bound the cut lines of one spanning tree $T$ of it by the first bound and the others by $C$.

\emph{Step 3: factorisation.} The output term is bounded by $\prod_K\mathfrak E_K\times\prod_{\ell\in T}Ce^{-\delta_2d(K,K')}\times C^{\#\text{cut lines}}$, where $\mathfrak E_K$ depends on the positions of $K$ only; the transfer factors cost one extra decay draw per line of a cluster, within the budget of \cref{thm:a02:WRall}.

\emph{Step 4: the cluster factors.} \emph{Zone-I clusters.} The cluster graph is a free $f_D$-flow graph with its cut lines as undifferentiated external legs of unit scale (\cref{lem:a02:XL}(b)). For the cluster containing the $h$-vertex, the census is that of \cref{prop:k02:KhI}; since the cluster carries a cut leg, every component has an external leg and the excepted field-independent family does not occur. For the source-free zone-I clusters, the degree count holds at every tuple by \cref{thm:a02:T} and \cref{thm:a02:WRall}, and block points by \cref{lem:a02:BLK}. The multiscale bound \cref{thm:a02:S3}(a) with the root held, summed over the root in $\Omega_K$ (volume $\le Cp^3$ by \cref{lem:a02:CL}(c)), gives $\sum_{\text{positions of }K}\mathfrak E_K\le C\eta^{\rho_K}p^{n_K}(1+k-j)^{p_{0,K}}$, with $\rho_K$ the coupling power of the vertices of $K$. \emph{Zone-B clusters.} The cluster is source-free (\cref{lem:k02:Zh}(d)); apply \cref{thm:a02:Bsc}(a) and \cref{thm:a02:BVP} with no vertex held to the renormalised family of $H_K$, localised at $P_K$ with the box $\Omega_K$ of side $O(p)$; the hypothesis ``at least one external field'' is the cut leg. The proof of \cref{thm:a02:Bsc} uses an external scalar leg only in the following ways, each of which a cut line satisfies: through its sup and $(1,\alpha)$-norms on $\Omega_K$ up to the faces, in the Taylor and transport steps, the H\"older splits and the ordered induction (\cref{lem:a02:XL}(b)); and, when the charging rule of \cref{lem:a02:CH} charges the leg at a $D'=-1$ remainder vertex $v$, through the vanishing factor $m_\alpha(v;1)$ of \cref{lem:a02:H}(ii), which a cut line supplies by \cref{lem:k02:face}(a). A cut line joining two zone-B clusters may be charged at both ends, by the charging rules of the two clusters; then \cref{lem:k02:face}(b) supplies both factors. After the borrowing of \cref{lem:a02:CH} the cut-line factor left in the output term is $\ell/[m_\beta(u';1)m_\beta(v';1)]$ (or with one factor), an undifferentiated evaluation bounded by $C\min(1,t(u'))^{1-\beta}\min(1,t(v'))^{1-\beta}e^{-\delta\dist}\le Ce^{-\delta\dist}$ by \cref{lem:k02:face}, which Steps~2--3 treat like every other cut-line factor. The resulting bound is again $\le C\eta^{\rho_K}p^{n_K}(1+k-j)^{p_{0,K}}$, the logarithm coming from the infrared sums inside $\Omega_K$.

\emph{Step 5: the tree.} By \cref{lem:a02:CL}(d), $\prod_{\ell\in T}e^{-\delta_2d(K,K')}\le e^{\delta_2c_0N}e^{-\delta_2\dist(P(\tau))}$.

\emph{Step 6: exponents.} $\rho(H)=\sum_K\rho_K$, the count being vertex-local (an unlocalised mass-counterterm vertex on a cut dressed line keeps its own coupling power); $n_1\defeq\sum_Kn_K$ and $p_0\defeq\sum_Kp_{0,K}$ plus the dressed-line factor of \cref{lem:a02:XL}(b); the number of clusters is $\le N_*$ and $\delta\defeq\delta_2$. This is \eqref{eq:k02:R1}.

\emph{The root sum.} On the tuples of \eqref{eq:k02:R2}, $R(\tau)\ge\eta^{-1}/2$ by \cref{lem:k02:Zh}(b). The number of tuples of $H$ with a point in $y_*$ and $R\in[n,n+1)$ is $\le C(n+1)^{3N_*}$, and $e^{\delta'R}e^{-\delta R}\le e^{-\delta R/2}$. So the sum is $\le C(1+k-j)^{p_0}\eta^{\rho(H)}p^{n_1}\sum_{n\ge\eta^{-1}/2-1}(n+1)^{3N_*}e^{-\delta n/2}\le C'(1+k-j)^{p_0}\eta^{\rho(H)}p^{n_1}e^{-\delta/(8\eta)}$, since $(n+1)^{3N_*}e^{-\delta n/4}\le C$ and $\eta^{-1}\ge8$. Members with curvature vertices: the curvature vertex lies in a zone-I cluster (\cref{lem:k02:Zh}(d)), is invariant under the reduction (\cref{lem:a02:SBP}(b)) and is bounded through \cref{lem:a02:SBP}(c).
\end{proof}

\begin{proof}[Proof of \cref{hyp:k02:Kh}]
Let $H$ be a member of $\mathcal N_h$ and $J_*\le j\le k-J_*$, so that $\eta\le\eta_2$. \emph{First clause.} \eqref{eq:k02:I1} is \cite[(3.39)]{K1} for the localisations in zone~I, which are those at distance from the faces larger than the small-field scale. \emph{Second clause.} A tuple with one cluster has $R\le c_0(N_*-1)^2\le p$, lies in zone~I by \cref{lem:a02:Z}(i), and obeys \eqref{eq:k02:I1}; a tuple with at least two clusters obeys \eqref{eq:k02:R1}. Summed over the localisations sharing the cube $y_*$ of an $h$-vertex: the zone-I tuples by $\sum_n(n+1)^{3N_*}e^{-(\delta-\delta')n}\le C$, the zone-B tuples by \eqref{eq:k02:R2}. This is \cite[(3.38)]{K1} with the factor $(1+k-j)^{p_0}$. \emph{Third clause.} Split the localisations at $R=p$. Those with $R\le p$ are in zone~I; sum \eqref{eq:k02:I2} by \cref{lem:a02:SC} (branch $R\le p$): the result is $\le C\eta^{\rho(H)-1/2}p^{n_1}(1+k)^{p_0}\eta^{-3}|\CJ|[5^{-\gamma j}+\eta^{3+\sigma+1/2}]$, whose first term is $\le C5^{-\gamma'j}\eta^{-\beta'}|\CJ|$ for every $\gamma'<\gamma$ and $\beta'>\frac72-\rho(H)$, the powers $p^{n_1}(1+k)^{p_0}$ being absorbed on the diagonal sequence as in \cref{lem:a02:SC}, and whose second term is $\le C\eta^{\rho(H)+\sigma}p^{n_1}(1+k)^{p_0}|\CJ|\le C\eta^{\sigma'}|\CJ|$ for every $\sigma'<\sigma$. Those with $R>p$ have at least two clusters; bound the two members of the pair separately by \eqref{eq:k02:R1}, with the roots at the $h$-vertex, $\le C\eta^{-3}|\CJ|$ of them. For the zone-I tuples among them the factor $e^{-(\delta-\delta')p}$ of \cref{lem:a02:SC} (branch $R>p$) is $\le C\eta^{3+\sigma}$, since $p(\eta)$ grows faster than every logarithm, and the sum is $\le C\eta^{-3}|\CJ|\eta^{\rho(H)}p^{n_1}(1+k-j)^{p_0}\eta^{3+\sigma}\le C\eta^{\sigma'}|\CJ|$; for the zone-B tuples \eqref{eq:k02:R2} gives $\le C\eta^{-3}|\CJ|\eta^{\rho(H)}p^{n_1}(1+k-j)^{p_0}e^{-\delta/(8\eta)}\le C_N\eta^N|\CJ|$ for every $N$. In both cases $\rho(H)\ge0$, and the external fields on $\supp\chi_j$ are accounted for in $\rho(H)$ by \eqref{eq:k02:margin}. Adding the three parts gives the third clause with $C$ independent of $j$ and $k$. The constants depend on the point of King's range, so the lemma holds at every point of the range, in particular at the second mass $M_2$ of \cref{ssec:k02:setting}.
\end{proof}



\subsection{The inputs with scalar sources}\label{ssec:k02:inputs}

For a source $s=(u,h)$ put $\mathsf S(s)\defeq\sum_i|x_i|\max_{\mu\nu}\norm{f_{i,\mu\nu}}_{C^{2,\alpha}}+\norm h_\infty$, a size; a compact family of sources is contained in $\{\mathsf S(s)\le S_0\}$.

\begin{proposition}[\Uplus: the small-field lines at a non-negative scalar source]\label{prop:k02:Uplus}
Fix $J\ge J_0$, $\alpha\in(0,\frac12)$ and $S_0>0$. There are $k_0$ and $C,\sigma>0$, depending on the couplings, $J$, $\alpha$, $S_0$ and on nothing else, such that for every real source $s=(u,h)$ with $h\ge0$ and $\mathsf S(s)\le S_0$, every $k\ge k_0$ and both members $(k',j')\in\{(k,j(k)),(k+1,j(k)+1)\}$ of the diagonal pair,
\begin{equation}\label{eq:k02:SF}
Z_{k'}(s)\ \ge\ e^{-C\eta_k^{\sigma}}\,I^{(j')}_{k'}(s)\qquad\text{and}\qquad \int P\,d\mu_{k',s}\ \le\ e^{C\eta_k^{\sigma}}\,I^{(j')}_{k'}(s).
\end{equation}
\end{proposition}
The proof is given in Appendix~\ref{app:a03_Uplus}. These are the two small-field lines of \eqref{eq:k01:H0rel} with the scalar source inserted into the small-field analysis through the vertex \cite[(3.21)]{K1}; the large-field remainder at $s\ne0$ is not part of the statement, and is supplied instead by \cref{lem:k02:HLF}. The source enters the small-field steps only through its vertex, and the sign $h\ge0$ is used only through the sign of the Wick-ordered source term in the large-field region, which is not examined here.

\begin{proposition}[\Bh: one-step bound with the scheme $c$-number]\label{prop:k02:Bh}
Fix $J\ge J_0$ and a compact family $\Sigma$ of real sources. For $\sigma\in\Sigma\cup\{0\}$ and $h_\sigma$ its scalar part, there is a sequence $\theta_k(\sigma)\le\theta_k$, with $\theta_k$ independent of $\sigma\in\Sigma$ and decaying geometrically, such that for $k\ge k_0$, on the support of $\chi_{j(k)}$,
\begin{equation}\label{eq:k02:B}
\big|S^{(j(k))}_{N,k}(\cdot;\sigma)+c_k(h_\sigma)-S^{(j(k)+1)}_{N,k+1}(\cdot;\sigma)-c_{k+1}(h_\sigma)\big|\ \le\ \theta_k(\sigma).
\end{equation}
\end{proposition}
The proof is given in Appendix~\ref{app:a04_Bh}. The field-dependent and the $h$-nonlinear parts are \cref{prop:k01:Kpos}, the graphs with curvature vertices at $s=(u,h)$, and \Kh. The field-independent $h$-linear part of the three lowest orders is a $c$-number $\bar\beta^{(j)}_k\int h$ up to a localisation error, and $c_k(h)$ is exactly the $c$-number that makes $\bar\beta^{(j)}_k+z_k$ converge along the diagonal pairs at the rate $5^{-\gamma'j}\eta^{-\beta'}+\eta^{\sigma'}$; this is the sense in which $z$ is the scheme adapted to King's flow. At $\sigma=0$, \eqref{eq:k02:B} is \eqref{eq:k01:Kpos} with $j=j(k)$.

\begin{lemma}[$c$-numbers]\label{lem:k02:cnum}
\textup{(a)} $|z_k|\le C(1+k)$. \textup{(b)} For every $M>0$, $\Delta_k(M)\defeq2(\sW_{\eps_k}(m)-\sW_{\eps_k}(M))+z_k(m)-z_k(M)$ is bounded in $k$, where $\sW_\eps(M)=(-\Delta_{\eps\Z^3}+M^2)^{-1}(0,0)$ and $z_k(M)$ is the reference scheme at mass $M$.
\end{lemma}
The proof is given in Appendix~\ref{app:a05_cnum_NEG}. Both are statements about lattice Feynman integrals: (a) is the logarithmic growth of the order-$e^2$ coefficient of $\langle|\phi|^2\rangle$ with the counterterm inserted; in (b) the Wick difference is a convergent lattice integral in three dimensions, and the scheme difference converges.

\begin{proposition}[\NEGinf: the bare coefficient grows with the mass]\label{prop:k02:NEG}
For every point of King's range ($e,\lambda\ge0$, $m,\lmu>0$) and every $M>m$, $u_\infty(M)\defeq\lim_ku_k(M)$ exists and
\[
2u_\infty(M)\ \ge\ (M-m)\Big(M+m+\frac{4\lambda}{\pi}-\frac{e^2}{4\pi}\Big)-\Big(\frac{e^4}{4\pi^2}+\frac{8\lambda^2}{\pi^2}\Big)\ln\frac Mm ;
\]
in particular $s_-(M)=u_\infty(M)\to\infty$ as $M\to\infty$, and $s_-(M_2)\ge\frac12$ for $M_2\defeq\max\{m,\ \frac{e^2}{4\pi}+(\frac{e^4}{4\pi^2}+\frac{8\lambda^2}{\pi^2})/m\}+1$.
\end{proposition}
The proof is given in Appendix~\ref{app:a05_cnum_NEG}. The content is the exact continuum limit of King's free-boundary counterterm at all five of its orders and the $k$-uniform bounds on its derivatives in the mass; only $\sup_Ms_-(M)=\infty$ is used below.

\begin{proposition}[{\Wsh$[s_W]$: the positive-source bound}]\label{prop:k02:Wsharp}
Assume \Itwo\ and \Ithree\ (\cref{prop:s10:U,prop:s10:RP}; \Itwo\ follows from \Tall). Then with $s_W\defeq s_1$ (the amplitude of \cref{prop:s11:W1prime}) and $\kappa=C_U/s_1$: for every $k$, $J$ and every lattice function $0\le q\le s_W$ on $\XJ$,
\[
\ln n_k(0,q)\ \le\ \kappa\,\norm q_{1,k}.
\]
In particular $\sup_kn_k(0,q)<\infty$ for every $0\le q\in C^\infty_c$ with $\sup q\le s_W$ and $\supp q$ at distance $>1$ from $\R^3\setminus\CJ$.
\end{proposition}
\begin{proof}
$e^{\sum_x\eps^3q(x)r_x^2}$ is a class-$\cI$ observable, so by Step~1 of \cref{lem:s04:G} (killing the matter on $\TJ$ does not increase it; \cref{rem:s04:F}) $n_k(0,q)\le N^z_{k,\CJ}(0,q)$, the comparison-cube functional, which \cref{prop:s11:W1prime} at $X=\XJ$ bounds by $e^{\kappa\norm q_{1,k}}$; the scheme multiplier is common (\cref{lem:s02:alg}(a)). For $q\in C^\infty_c$, $\norm q_{1,k}\le C_0'\wnorm q$ by \cref{lem:s02:E}(E3).
\end{proof}
Any other threshold with the same property serves below in place of $s_1$; the appendix contains a second instance, obtained from Ba{\l}aban's torus bounds through the comparison of the cube with King's tori.

\subsection{Three lemmas on positive and negative scalar sources}\label{ssec:k02:lemmas}

For a mass $M$ let $Y^{(M)}_k(u,h)\defeq e^{-c_k(h)}\langle e^{F(u)+\Wick{h}}\rangle^{(M)}$ denote the functional of the mass-$M$ state \emph{in the conventions of the mass $m$} (free Wick constant $\sW_{\eps_k}(m)$ and scheme $z_k=z_k(m)$), so that $Y^{(m)}_k=n_k$ and $Y^{(M)}_k(u,h)=e^{-\Delta_k(M)\sS_k(h)}n^{(M)}_k(u,h)$, with $n^{(M)}_k$ the functional in the mass-$M$ conventions and $\Delta_k(M)$ as in \cref{lem:k02:cnum}(b). Write $r^2(q)\defeq\sum_{x}\eps_k^3q(x)r_x^2$ over the live sites.

\begin{lemma}[mass shift and the correlation inequality]\label{lem:k02:GM}
Let $M>0$, $k\ge0$, and let $0\le q$ be a lattice function on the live sites with $q(x)\le u_k(M)$ for all $x$. Then
\[
n_k(0,-q)\ \le\ \frac1{Y^{(M)}_k(0,q)},\qquad\text{and}\qquad \frac1{n_k(0,q)}\ \le\ n_k(0,-q)\ \le\ \frac1{Y^{(M)}_k(0,q)},\quad\text{hence}\quad Y^{(M)}_k(0,q)\le n_k(0,q).
\]
\end{lemma}
\begin{proof}
\emph{Mass shift.} In King's action the scalar mass enters only through $\sum_x\eps^3\frac12(M^2+\dm_\eps(M))r_x^2$: $\dm_\eps$ is one number per spacing and mass (\cref{def:s02:sources}), the Dirichlet terms $\frac\eps2r_x^2$ do not depend on $M$, and the photon does not see $M$. Hence, with $u\defeq u_k(M)$ and $r^2(1)\defeq\sum_x\eps_k^3r_x^2$, $e^{-S^{(m)}}=e^{-S^{(M)}}e^{u\,r^2(1)}$, and for every $\Phi$ integrable against both measures $\langle\Phi\rangle^{(m)}=\langle\Phi e^{ur^2(1)}\rangle^{(M)}/\langle e^{ur^2(1)}\rangle^{(M)}$; the quartic term makes $e^{\pm r^2(\cdot)}$ integrable. \emph{Correlation inequality.} Put $g_1\defeq u\mathbf 1-q\ge0$ and $g_2\defeq q\ge0$. The mass-$M$ state is a BFS measure, and so is the state with its single-site measures multiplied by $e^{t\,r^2(g_1)}$, $t\in[0,1]$ (\cref{def:s02:acube} and the paragraph on BFS measures following it: the exponential moments stay finite because $\lambda>0$). By \cref{lem:s02:BFS}(c), multiplying the single-site measures by $e^{-tr^2(g_1)}$ does not increase $\langle e^{r^2(g_2)}\rangle$; read from $t=1$ down to $t=0$, $\langle e^{r^2(g_2)}\rangle^{(M)}_{t=1}\ge\langle e^{r^2(g_2)}\rangle^{(M)}_{t=0}$, that is
\[
\langle e^{r^2(g_1)}e^{r^2(g_2)}\rangle^{(M)}\ \ge\ \langle e^{r^2(g_1)}\rangle^{(M)}\langle e^{r^2(g_2)}\rangle^{(M)}.
\]
\emph{Conclusion.} With $\Phi=e^{-r^2(q)}$: $\langle e^{-r^2(q)}\rangle^{(m)}=\langle e^{r^2(g_1)}\rangle^{(M)}/\langle e^{r^2(g_1+g_2)}\rangle^{(M)}\le1/\langle e^{r^2(q)}\rangle^{(M)}$. Multiply both sides by $e^{(2\sW_{\eps_k}(m)+z_k)\sS_k(q)}$: the left side becomes $n_k(0,-q)$ and the right side $1/Y^{(M)}_k(0,q)$, both in the conventions of the mass $m$. The left inequality of the second display is \cref{lem:s02:alg}(e), and the last inequality follows.
\end{proof}

\begin{lemma}[the large-field remainder at a source, from the one at zero source]\label{lem:k02:HLF}
Let $\sigma$ be a real source and $p>1$, and assume \eqref{eq:k02:SF} at $\sigma$ and \Rzero\ of \cref{prop:k01:H0rel} for the member $(k,j)$, with $a\defeq C\eta^\sigma$ and $\rho\defeq\rho(\eta)$. Then
\[
e^{-a}\,\Ik{j}(\sigma)\ \le\ \Zk(\sigma)\ \le\ e^{a}\,\Ik{j}(\sigma)+\rho^{1-1/p}\,\langle e^{X(p\sigma)}\rangle_k^{1/p}\,\Zk(0).
\]
\end{lemma}
\begin{proof}
The lower line is the first inequality of \eqref{eq:k02:SF}. For the upper line, $\Zk(\sigma)=\int P\,d\mu_{k,\sigma}+\int(1-P)e^{X(\sigma)}d\mu_{k,0}$. The first term is at most $e^{a}\Ik{j}(\sigma)$ by \eqref{eq:k02:SF}. For the second, Hölder's inequality for the positive measure $(1-P)d\mu_{k,0}$ with the functions $1$ and $e^{X(\sigma)}$, then $1-P\le1$ in the second factor and $pX(\sigma)=X(p\sigma)$ (linearity), give
\[
\int(1-P)e^{X(\sigma)}d\mu_{k,0}\ \le\ \Big[\int(1-P)d\mu_{k,0}\Big]^{1-1/p}\Big[\int e^{X(p\sigma)}d\mu_{k,0}\Big]^{1/p}\ \le\ \big[\rho\Zk(0)\big]^{1-1/p}\big[\Zk(0)\langle e^{X(p\sigma)}\rangle_k\big]^{1/p},
\]
using \Rzero\ for the first bracket.
\end{proof}
At a negative scalar source the same step would need a bound on $\langle e^{X(p\sigma)}\rangle$ at a negative source of larger amplitude, which is what is to be proved; \cref{thm:k02:UB} avoids this through the mass shift.

\begin{lemma}[positivity of the limit at a positive source]\label{lem:k02:POS}
Fix a mass $M$, a function $0\le q\in C^\infty_c$ as in \cref{prop:k02:Wsharp}, and $\bar w\in(0,1]$. Assume that at mass $M$, in its own conventions, the hypotheses of \cref{thm:k02:CAh} below hold at $\sigma=(0,wq)$ for every $w\in[0,\bar w]$ and at $\sigma=0$, with sequences $b_k$, $\theta_k(wq)$, $r_k(0)$ and $e^{-c_k(wq)}r_k(wq)$ dominated by $w$-independent summable sequences, for all $k\ge k_0$ with $k_0$ independent of $w$. Then $\inf_{k\ge k_0}n^{(M)}_k(0,q)>0$.
\end{lemma}
\begin{proof}
Steps~1--4 of the proof of \cref{thm:k02:CAh} use the hypotheses at $\{wq,0\}$ only and give the lower half of \eqref{eq:k02:rec} with $w$-independent summable $\bar\delta_k$, $\bar\epsilon_k$: $n_{k+1}(wq)\ge(1-\bar\delta_k)n_k(wq)-\bar\epsilon_k$. By induction, $n_k(wq)\ge\prod_{i=k_1}^{k-1}(1-\bar\delta_i)\,n_{k_1}(wq)-\sum_{i\ge k_1}\bar\epsilon_i$ for $k\ge k_1$. Choose $k_1\ge k_0$ with $\prod_{i\ge k_1}(1-\bar\delta_i)\ge\frac12$ and $\sum_{i\ge k_1}\bar\epsilon_i\le\frac18$. At this fixed $k_1$, $w\mapsto n_{k_1}(wq)$ is continuous (a finite-dimensional integral) and equals $1$ at $w=0$; choose $w_0\in(0,\bar w]$ with $n_{k_1}(w_0q)\ge\frac12$. Then $n_k(w_0q)\ge\frac18$ for all $k\ge k_1$. Since $\ln n_k(wq)$ is convex in $w$ and vanishes at $w=0$, $n_k(w_0q)\le n_k(q)^{w_0}$, so $n_k(q)\ge8^{-1/w_0}$ for $k\ge k_1$; the finitely many $k_0\le k<k_1$ have $n_k(q)>0$.
\end{proof}
